\documentclass{article}
\usepackage[utf8]{inputenc}
\usepackage[T1]{fontenc}
\usepackage{geometry}
\usepackage{graphicx}
\usepackage{float}
\usepackage{hyperref}
\usepackage[english,indonesian]{babel}

\title{Optimasi Pemrosesan Template LaTeX Menggunakan Pendekatan Berbasis AI}
\author{Academic Support Agent}
\date{\today}

\begin{document}
\maketitle

\begin{abstract}

\end{abstract}

\chapter{PENDAHULUAN}


\chapter{TINJAUAN PUSTAKA}


% Original content (commented out):
% % **********************************************************************************************
% **********************************************************************************************
% 
%   Master thesis template of D-HEST, ETH Zurich
% 
%   This is an example template for master thesis in D-HEST, ETH Zurich.
%   - Upload the figures and graphs directly to the same directory.
%   - Export your literate list from, for example, EndNote, to a text file and then copy and paste the list to the file named bibliography.bib.
%   - More tutorials are found in https://www.overleaf.com/learn/latex/Tutorials.
%   - Follow the D-HEST guideline for details: https://hest.ethz.ch/en/studies/health-sciences-and-technology/documents.html
% 
% **********************************************************************************************
% **********************************************************************************************



% ##############################################################################################
% 
%   Initial settings
% 
% ##############################################################################################
%   Set the document format such as font, style.
%   https://www.overleaf.com/learn/latex/Creating_a_document_in_LaTeX
\documentclass[12pt]{article}

%   Load packages.
\usepackage[a4paper, right=15mm, left=25mm, top=20mm, bottom=20mm]{geometry}
\usepackage{fancyhdr, lastpage, graphicx, wrapfig, subcaption, setspace, booktabs, url, lipsum, hyperref, sectsty, multirow, color, tabularx, bibentry, bbding, framed, natbib, fourier}
\usepackage[T1]{fontenc} % to use font of Times New Roman
\usepackage[protrusion=true, expansion=true]{microtype} % to improve the font style beautifully.
\usepackage[font=small, labelfont=bf]{caption} % to change the font style in caption of figures and tables.
\usepackage[UKenglish]{babel} % to set British English.
\usepackage[acronym]{glossaries} % to list glossaries.
\usepackage[utf8]{inputenc} % to output accented letters and some symbols directly.

%   Set the space of lines. (1.0 for single space, 1.3 for one and half space, 1.6 for double space)
%   https://www.overleaf.com/learn/latex/paragraph_formatting
\linespread{1.3}

%   Set the page style; headers and footers. (plain, fancy, etc.s)
%   https://www.overleaf.com/learn/latex/headers_and_footers
\pagestyle{plain}

%   Start to create a document.
\begin{document}


% ##############################################################################################
% 
%   Title page
% 
% ##############################################################################################
\begin{titlepage}
    \begin{center}
        
        \Large
        ETH Zurich                                      \\
        Department of Health Sciences and Technology    \\
        Master in Health Sciences and Technology        \\
        Major in \textcolor{red}{[xxx]}
        
        \vspace{1.5cm}
 
        \framebox{
        \Huge       
        \textbf{\textcolor{red}{[TITLE]]}}
        }
        
        \vspace{1.5cm}
        \Large
        \textbf{Master Thesis}
        
        \vspace{1cm}
        \large
        In partial fulfilment of the requirements for the Master of Science ETH \\
        in Health Sciences and Technology (MSc ETH HST)                         \\
        
        \vspace{1cm}
        \large
        Submitted by:   \\
        \textcolor{red}{[Last name   First name]}
        
        \vspace{1cm}
        \large
        Supervisor(s):  \\
        \textcolor{red}{[(Title) Last name  First Name]}   \\
        \textcolor{red}{[Institute or Professorship or Institution]}   \\

        \vspace{1cm}
        \large
        Tutor:  \\
        \textcolor{red}{[(Title) Last name  First Name]}   \\
        \textcolor{red}{[Institute or Professorship or Institution]}   \\
        
        \vspace{1cm}
        \large
        Zurich, \textcolor{red}{[Month Year]}

        \vspace{1cm}
        \includegraphics[width=6cm]{ETH_logo_short.png}

    \end{center}
\end{titlepage}



% ##############################################################################################
% 
%   Abstract
% 
% ##############################################################################################
\section*{Abstract}
This is abstract part. If you do not want to assign numbers to sections, put * before the bracket.



% ##############################################################################################
% 
%   Show the contents of whole paper and summary of figures and tables. (no need to edit)
% 
% ##############################################################################################
\newpage
%   List the contents.
\tableofcontents
%   List up brief explanation of figures and tables.
\listoffigures
\listoftables

\clearpage



% ##############################################################################################
% 
%   1. Introduction
% 
% ##############################################################################################
\section{Introduction}
\subsection{Section name}
\subsubsection{Subsection name}
The following is the example to illustrate a figure. 

\subsection{Example: how to create figures?}
\textcolor{blue}{Example: Figure \ref{example_figure_1} shows the example image.}

When you put figures, you use the following "figure" tool. "centerling" sets the position of figure to centre. "includegraphics" calls the figure that you want to put (i.e. \verb|Figure_1-age_map.png| in this case). You can change the size of figure by adjusting the value of "width". "caption" is to explain what the figure shows. You can put any names to "label", which you call in the sentences.

Please refer to the following page for more tutorials: \url{https://www.overleaf.com/learn/latex/Inserting_Images}.


\begin{figure}[htbp]
    \centering
    \includegraphics[clip, width=1.0\textwidth]{Figure_1.jpg}
    \caption{Example image (\textit{Example})}
    \label{example_figure_1}
\end{figure}


\subsection{Example: how to cite literature?}
\textcolor{blue}{Example: It was estimated that the number of people living with dementia worldwide will increase from 35.6 million in 2010 to 65.7 million in 2030 and to 115.4 million in 2050, with almost double every 20 years \cite{RN293, RN189, RN181}.}

Literature data can be exported using, for example, EndNote BibTex format. All the literature data are just copied and pasted to "bibliography.bib". Then, you just need to call the reference numbers that are automatically allocated on EndNote, by using "citep" or "cite".

Please refer to the following page for more tutorials: \url{https://www.overleaf.com/learn/latex/Bibliography_management_in_LaTeX}.


\subsection{Example: how to create a table?}
\textcolor{blue}{Example: Table \ref{demographic} shows the demographic information on participants.}

When you create tables, you use the following "table" tool. Please refer to the page below for more tutorials: 
\url{https://www.overleaf.com/learn/latex/tables}.


\begin{table}[htbp]
    \setlength\arrayrulewidth{1pt}
    \centering
    \small
    \caption{Demographic profile of participants}
    
    %   Set the number of columns
    \begin{tabular}{lrl}
    %   Put a horizontal line.
    \hline
    \textbf{Variable} & \textbf{Values} & \textbf{P value}    \\
    \hline
    \hline
    Gender      & Men: 6, Women:    &               \\
    Age (years) & $73.14\pm4.26$    &   $P<0.001$   \\
    Weight (kg) & $77.26\pm19.39$   &   $P=0.201$   \\
    Height (cm) & $168.43\pm7.91$   &   $P=0.444$   \\
    BMI         & $27.00\pm5.21$    &   $P=0.012*$  \\
    \hline
    \end{tabular}
    \label{demographic}
\end{table}



\subsection{Ignore specific parts from counting the number of words.}
You can eliminate some parts from counting the total number of words, by using TC: ignore.
Please refer to the following page for more tutorials: \url{https://www.overleaf.com/learn/how-to/Is_there_a_way_to_run_a_word_count_that_doesn%27t_include_LaTeX_commands%3F}.

%%TC:ignore 
This part is not counted for the total number of words.
%%TC:endignore



% ##############################################################################################
% 
%   2. Research Method
% 
% ##############################################################################################
\section{Research Method}
\subsection{Experimental system}
\subsection{Measurement protocol}
\subsection{Data processing}
\subsection{Statistical analysis}



% ##############################################################################################
% 
%   3. Results
% 
% ##############################################################################################
\section{Result}
\subsection{Demographic data}


% ##############################################################################################
% 
%   4. Discussion
% 
% ##############################################################################################
\section{Discussion}
\subsection{Interpretation of results}
\subsection{Limitations and improvements}


% ##############################################################################################
% 
%   5. Conclusion
% 
% ##############################################################################################
\section{Conclusion}



% ##############################################################################################
% 
%   Acknowledgements
% 
% ##############################################################################################
\section*{Acknowledgements}
Example: I would like to thank the whole team for supporting this research project.



% ##############################################################################################
% 
%   Reference
%   Need to edit bibliography.bib. (Export literature data from literature management 
%   tool such as EndNote. EndNote provides export style called "BibTeX".)
% 
% ##############################################################################################
\newpage
\bibliography{bibliography}

%   Set the style to show the literature.
%   https://www.overleaf.com/learn/latex/bibtex_bibliography_styles
% \bibliographystyle{apalike}
\bibliographystyle{unsrt}


%   End the document.
\end{document}



\end{document}